%%%%%%%%%%%%%%%%%%%%%%%%%%%%%%%%%%%%%%%%%
% Laboratory Report LaTeX Template
%
% This template has been downloaded from:
% http://www.latextemplates.com
%
%%%%%%%%%%%%%%%%%%%%%%%%%%%%%%%%%%%%%%%%%

%----------------------------------------------------------------------------------------
%	DOCUMENT CONFIGURATIONS
%----------------------------------------------------------------------------------------

\documentclass{article}

\usepackage{wasysym}
\usepackage{amssymb}% http://ctan.org/pkg/amssymb
\usepackage{pifont}% http://ctan.org/pkg/pifont
\newcommand{\cmark}{\ding{51}}
\newcommand{\xmark}{\ding{55}}

\usepackage{graphicx}

\title{Apparent Positions of The Planets \\ ASTR 101} % Title

\author{Jordan \textsc{Ell} \\ jordan.ell7@gmail.com \\ V00660306} % Author name

\begin{document}

\maketitle % Insert the title, author and date

\begin{tabular}{lr}
Date Performed: 24/09/2012 \\ % Date the experiment was performed and partner's name
Supervisor: Jillian Scudder % Instructor/supervisor
\end{tabular}

\setlength\parindent{0pt} % Removes all indentation from paragraphs

\renewcommand{\labelenumi}{\alph{enumi}.} % Make numbering in the enumerate environment by letter rather than number (e.g. section 6)

%----------------------------------------------------------------------------------------
%	SECTION 1
%----------------------------------------------------------------------------------------

\section{Objective}

This report will demonstrate the proper procedure of finding planets apparent positions in the sky.
It explains how to map the planets on a polar coordinate graph and ecliptic constellation chart as well as the 
use of a planisphere.\\
 
%----------------------------------------------------------------------------------------
%	SECTION 2
%----------------------------------------------------------------------------------------

\section{Introduction}

Observation of the planets is an important task in the field of astronomy. The term planet comes from an Ancient Greek word
meaning "wandering star". Ancient tribes first spotted these stars moving against the backdrop of other stars and have been
since studied for their orbits, elemental make-up as well as other properties. While most planets in our solar system have been
discovered by simply looking up at the sky, some have also been discovered by mathematics (Neptune)~\cite{Edmund:2008}. 
As science has evolved, our model for the planets has moved from an Earth to Sun centred, which has drastic effects on the 
planet's orbits. We have refined our notion of planetary movement from around the Earth to perfect circles around the Sun,
to Kepler's 1st law which explains the elliptical path that planets take around the Sun~\cite{Holton:2009}. Our studies of the
planets more recently have even included multiple satellites and even a few rovers on Mars to search for evidence of past
life~\cite{NASA:2011}. While these explorations are viable, this report will focus on a simpler understanding of some planetary
properties.
After having read this report, our view of how to find the planets in the sky should become clear as well as having a broad 
understanding of their positions in both distance and period from the Sun, and their equatorial positioning from Earth.
During this lab, we will make the large assumption that planets travel in perfect circles around the Sun. We know this not
to be true but will make for easier observations and possible calculations. It will also be assumed that planets travel along
the ecliptic in a perfect path.\\

%----------------------------------------------------------------------------------------
%	SECTION 3
%----------------------------------------------------------------------------------------

\section{Equipment}

The three main pieces of equipment required for this lab are as follows. First, a polar coordinate graph which is used for
charting where planets are relative to the sun at any given date. This report will use the chart with measuring 1 AU as 
one dark black ring from the centre. Second, a  SC001 constellation equatorial chart. This chart has measurements in declination, 
measured in degrees, and right ascension, measured in hours and minutes. Lastly, a planisphere is needed, which is a
quick reference to stars given a day and time of night. Other common equipment needed for this lab includes: protractor,
scissors, and coloured pencils.\\


%----------------------------------------------------------------------------------------
%	SECTION 4
%----------------------------------------------------------------------------------------

\section{Procedure}

Given the data inside of Table~\ref{tab:long} and Table~\ref{tab:radii}, a heliocentric sketch was made of Venus', Earth's, Mars', 
Jupiter's, and Saturn's position at the dates of 2012-Sept-22, 2012-No-21, and 2013-Mar-20. This was accomplished by first
looking up a planet in Table~\ref{tab:radii} to find its orbit radius (measured in A.U) in order to give it a distance away from the
sun. We are assuming perfect circles for orbits. Next, a lookup of the planet's heliocentric longitude was preformed on 
Table~\ref{tab:long}. With these two coordinates, a sketch of where the planet should be located was able to be made on the 
polar coordinate graph given that the Sun was placed at the centre and that each dark black line represented 1 A.U. This process
was repeated for all five planets and all three dates listed in Table~\ref{tab:long}\\

Next, the visibility of each planet given the four times of day: noon, sunset, midnight, and sunrise are made. This was accomplished
using a small piece of paper as a guide for East and West on the Earth placed over top of the Earth in the sketch made above. Knowing,
the Earth rotates counter clockwise, if viewed from above, the paper guide was rotated to simulate a field of vision on the Earth at
the four given times of day. A planet's visibility can then be determined by locating it (visible) or not locating it (not visible) in the
180$^\circ$ view provided by the paper guide. These results can be found in Table~\ref{tab:tod}.\\

Now, the determination of the planet's alignment is made. By looking at sketch made above, it was determined (for the day of
2012-Sept-22) what a planet's alignment is relative to the Earth and Sun. This was done purely by observation of the sketch
made above and with 
approximating the planet's actual position to that of a specific alignment. The results from these observations can be found
in the Configuration column of Table~\ref{tab:tod}.\\

With the now sketched heliocentric view of planetary positions, geocentric equatorial positions of the planets were then
calculated for the day of 2012-Sept-22. This was accomplished by first measuring the geocentric ecliptic longitudes of the planets.
This was calculated by measuring the angle between the Earth and the planets in question, noting that the Earth is now the origin 
of measurement as opposed to the Sun that angles are measured from. These measurements can be found in the 
Ecliptic Long. column of Table~\ref{tab:pos}. Next, a plot of the ecliptic longitudes was made onto the provided SC001
constellation chart. This was accomplished by simply mapping the found ecliptic longitude above with the given degrees on
the ecliptic on the SC001. This is under the assumption that the planets in question follow the ecliptic path perfectly and do not
stray from it. Finally, the right ascension (measured in hours and minutes) and declination (measured in degrees)
were measured off the SC001 chart for the Sun and given planets. The results of these readings can be found in Table~\ref{tab:pos}.
These results are our final geocentric equatorial positions of the planets for the date of 2012-Sept-22.\\

Lastly, minor approximate measurements were made using a Planisphere. The Planisphere was used by placing a provided circular
map of the stars with labelled months and days of the year underneath an ellipse which has times of night and rotating the 
circular map of the stars so that your desired time and day match up. What is shown through the ellipse of the stars is what will 
be visible from the given latitude of the planisphere. The questions and answer measurements taken can be found in
Section~\ref{sec:qna}.\\


%----------------------------------------------------------------------------------------
%	SECTION 5
%----------------------------------------------------------------------------------------

\section{Observations}

All observations for this report were made on the day of 24/09/2012. As this report is based off of given numerical data as
well as paper observations, weather conditions and time of day have not been reported as they hold no effect. For all 
observations made using the polar coordinate chart, the chart and its sketches can be found as Figure 1 in the back of
this lab report. These observations include: the planet's position at the given dates in Table~\ref{tab:long}, ecliptic 
longitude of planets, and approximations of planet's alignment. All observations made using the SC001 constellation
chart can be found as Figure 2 in the back of this report. These observations include both the measurements of right
ascension and declination of the sun and planets.\\

%----------------------------------------------------------------------------------------
%	SECTION 6
%----------------------------------------------------------------------------------------

\section{Tables and Measurements}
\label{sec:tnm}

\begin{table}[h!]
\begin{center}
\begin{tabular}{l c c c c c }
\hline
Date & Venus & Earth & Mars & Jupiter & Saturn\\
\hline
\hline
2012-Sep-22 & 068 & 000 & 261 & 065 & 211\\
2012-Nov-21 & 166 & 060 & 297 & 070 & 213\\
2013-Mar-20 & 354 & 180 & 011 & 081 & 217\\
\hline
\end{tabular}
\end{center}
\caption{Heliocentric Longitudes from Astronomical Almanac.\label{tab:long}}
\end{table}

\begin{table}[h!]
\begin{center}
\begin{tabular}{ c c c c }
\hline
Planet & Orbit Radius (A.U) & Period (years) & Symbol\\
\hline
\hline
Sun &  &  & \astrosun \\
Venus & 0.72 & 0.62 & \venus\\
Earth & 1.00 & 1.00 & \earth\\
Mars & 1.52 & 1.88 & \mars\\
Jupiter & 5.20 & 11.86 & \jupiter\\
Saturn & 9.54 & 29.46 & \saturn\\
\hline
\end{tabular}
\end{center}
\caption{Radii and Period of Orbits.\label{tab:radii}}
\end{table}

\begin{table}[h!]
\begin{center}
\begin{tabular}{| c | c | c | c | c | c |}
\hline
Planet & Noon & Sunset & Midnight & Sunrise & Configuration\\
\hline
\hline
Venus & \cmark & & & \cmark & Greatest Western Elongation \\
Mars & \cmark & \cmark & & & Greatest Eastern Elongation\\
Jupiter & & & \cmark & \cmark & Quadrature\\
Saturn & \cmark & \cmark & & & Conjunction \\
\hline
\end{tabular}
\end{center}
\caption{Planets as seen from the Earth.\label{tab:tod}}
\end{table}

\begin{table}[h!]
\begin{center}
\begin{tabular}{| c | c | c | c | c | }
\hline
Planet & Ecliptic Long. & Constellation & Right Ascension & Declination\\
\hline
\hline
Sun & 180$^\circ$ & Virgo & 12:00 & 0$^\circ$. \\
Venus & 134$^\circ$ & Cancer & 09:02 & 16$^\circ$. \\
Mars & 230$^\circ$ & Libra & 15:10 & -17$^\circ$. \\
Jupiter & 73.5$^\circ$ & Taurus & 04:41 & 21$^\circ$. \\
Saturn & 207$^\circ$ & Virgo & 13:42 & -11$^\circ$. \\
\hline
\end{tabular}
\end{center}
\caption{Geocentric Equatorial Position of the Planets.\label{tab:pos}}
\end{table}

%----------------------------------------------------------------------------------------
%	SECTION 7
%----------------------------------------------------------------------------------------

\section{Questions}
\label{sec:qna}

The following questions and answers are asked throughout the lab manual for ASTR101 Lab 2. The
questions have been repeated for the reader.

\begin{enumerate}
\item[Q.] Why does Venus go through phases? 
\item[A.] Venus has phases similar to that of the moon because of the view the Earth has of Venus as it travels
around the Sun \textit{inside} of the orbit of the Earth around the Sun. The different phases are caused by the different
perspective we have on Earth of the sun shining on Venus.  The first person to observe these phases was Galileo 
Galilei in 1610~\cite{Lance:1995}. Figure~\ref{fig:pov} below shows how the different phases of
Venus appear from Earth.

\addtocounter{figure}{2}
\begin{figure}[h!]
\centering
\includegraphics[width=0.4\textwidth]{images/blank}
\caption{The phases of Venus.\label{fig:pov}}
\end{figure}

\item[Q.] In which zodiacal constellation was the Sun located when you were born? What is your astrological sign?
Check with your partners and discuss any discrepancy.
\item[A.] As per Figure 2, the sun was located in Gemini when I was born on July 12th. This is a discrepancy with
my astrological sign of Cancer. This difference is caused by the precession of the Earth. Unknown to most, there
is actually a 13th constellation which the Sun passes through known as Ophiuchus as opposed to the 
traditionally known 12~\cite{Storm:1995}.
\end{enumerate}

The following answers are labelled to coincide with the lab manual for ASTR101 Lab 2 section
"The Use of a Planisphere". As
questions 1 and 3 are instructions rather than actual questions/measurements, they have been skipped. All answers 
are given using the provided planisphere and approximations.

\begin{enumerate}
\item[2.] While the circular dial is being rotated, the star Polaris does not seem to move. This is to be expected as
Polaris the "North Star" is the \textit{current} pivot point around which all stars rotate.
\item[4.] At the time of 12AM on the night of June 01, the star Vega, which is apart of the constellation Lyra, appears to 
be in the centre (Zenith) of the planisphere. Vega was the northern pole star around the year 12,000 BC~\cite{Stell:2009}.
\item[5.] At the time of 12AM on the night of June 01, the star Spica, which is apart of the constellation Virgo, appears
to be on the "Western Horizon" of the planisphere.
\item[6.] At the time of 11PM on the night of June 15, the star Vega again appears at the Zenith while the star Spica
again appears at the "Western Horizon" of the planisphere. Spica was a key star in the discovery of the precession
of the Earth~\cite{Rufus:1943}.
\item[7.] At the time of 11PM on the night of August 15, the star Deneb, which is apart of the constellation Cygnus, appears
to be at the Zenith, while the star Arcturus appears to be at the "Western Horizon" of the planisphere.
\item[8.] On the night of September 22nd, the star Antares will set at approximately 9PM. The star Antares will also
set at 1AM on the night of July 30th.
\item[9.] On the night of January 1st, the star Sirius will rise at approximately 7PM and will set at approximately 5AM, 
giving a total time of 10 hours above the horizon. Sirius is the brightest star in the night sky partial due to its proximity
to Earth at a distance of 2.6 parsecs~\cite{van:2007}.
\end{enumerate}

%----------------------------------------------------------------------------------------
%	SECTION 8
%----------------------------------------------------------------------------------------

\section{Conclusions/Discussions}

This report has shown how to locate a planet in the sky given initial data points for orbit radius and heliocentric 
longitudes. This report has also shown the use of a planisphere for quick star reference at a given latitude.

%----------------------------------------------------------------------------------------
%	SECTION 9
%----------------------------------------------------------------------------------------

\section{Evaluation}

I found this report very useful for my own observational interest. I recently acquired a Celestron Super C8 Plus and have
been interested in observing the planets but have not known how to properly find them in the sky. I now have the 
required skills to be able to find planets and make predictions of their whereabouts throughout the year. However, I would
have liked to have had the opportunity to actually take a measurement of the planet rather than just have the initial
data provided.

%----------------------------------------------------------------------------------------
%	BIBLIOGRAPHY
%----------------------------------------------------------------------------------------

\begin{thebibliography}{9}

\bibitem{Edmund:2008}
O'Connor, John J.; Robertson, Edmund F. (2006). "John Couch Adams' account of the discovery of Neptune". University of St Andrews. Retrieved February 18, 2008.

\bibitem{Holton:2009}
Holton, Gerald James; Brush, Stephen G. (2001). Physics, the Human Adventure: From Copernicus to Einstein and Beyond (3rd paperback ed.). Piscataway, NJ: Rutgers University Press. pp. 40–41. ISBN 0-8135-2908-5. Retrieved December 27, 2009.

\bibitem{NASA:2011}
NASA Staff (26 November 2011). "Mars Science Laboratory". NASA. Retrieved 2011-11-26.

\bibitem{Lance:1995}
 Goines, David Lance (October 18, 1995). "INFERENTIAL EVIDENCE FOR THE PRE-TELESCOPIC SIGHTING OF THE CRESCENT VENUS". Retrieved 2010-08-11.

\bibitem{Storm:1995}
Kollerstrom, N. (October 1995). "Ophiuchus and the media". The Observatory (KNUDSEN; OBS) 115: 261–262.

\bibitem{Stell:2009}
Calculation by the Stellarium application version 0.10.2, retrieved 2009-07-28

\bibitem{Rufus:1943}
Rufus, W. Carl (April 1943). "Copernicus, Polish Astronomer, 1473–1543". Journal of the Royal Astronomical Society of Canada 37 (4): 134. Bibcode 1943JRASC..37..129R

\bibitem{van:2007}
van Leeuwen, F. (November 2007), "Validation of the new Hipparcos reduction", Astronomy and Astrophysics 474 (2): 653–664

\end{thebibliography}

\end{document}